\documentclass[10pt, letterpaper]{article}

% Packages:
\usepackage[
    ignoreheadfoot, % set margins without considering header and footer
    top=2 cm, % seperation between body and page edge from the top
    bottom=2 cm, % seperation between body and page edge from the bottom
    left=2 cm, % seperation between body and page edge from the left
    right=2 cm, % seperation between body and page edge from the right
    footskip=1.0 cm, % seperation between body and footer
    % showframe % for debugging 
]{geometry} % for adjusting page geometry
\usepackage{titlesec} % for customizing section titles
\usepackage{tabularx} % for making tables with fixed width columns
\usepackage{array} % tabularx requires this
\usepackage[dvipsnames]{xcolor} % for coloring text
\definecolor{primaryColor}{RGB}{0, 0, 0} % define primary color
\usepackage{enumitem} % for customizing lists
\usepackage{fontawesome5} % for using icons
\usepackage{amsmath} % for math
\usepackage[
    pdftitle={CV},
    pdfauthor={Ahmet Neçirvan Doğan},
    pdfcreator={Ahmet Neçirvan Doğan},
    colorlinks=true,
    urlcolor=primaryColor
]{hyperref} % for links, metadata and bookmarks
\usepackage[pscoord]{eso-pic} % for floating text on the page
\usepackage{calc} % for calculating lengths
\usepackage{bookmark} % for bookmarks
\usepackage{lastpage} % for getting the total number of pages
\usepackage{changepage} % for one column entries (adjustwidth environment)
\usepackage{paracol} % for two and three column entries
\usepackage{ifthen} % for conditional statements
\usepackage{needspace} % for avoiding page brake right after the section title
\usepackage{iftex} % check if engine is pdflatex, xetex or luatex

% Ensure that generate pdf is machine readable/ATS parsable:
\ifPDFTeX
    \input{glyphtounicode}
    \pdfgentounicode=1
    \usepackage[T1]{fontenc}
    \usepackage[utf8]{inputenc}
    \usepackage{lmodern}
\fi

\usepackage{charter}

% Some settings:
\raggedright
\AtBeginEnvironment{adjustwidth}{\partopsep0pt} % remove space before adjustwidth environment
\pagestyle{empty} % no header or footer
\setcounter{secnumdepth}{0} % no section numbering
\setlength{\parindent}{0pt} % no indentation
\setlength{\topskip}{0pt} % no top skip
\setlength{\columnsep}{0.15cm} % set column seperation
\pagenumbering{gobble} % no page numbering

\titleformat{\section}{\needspace{4\baselineskip}\bfseries\large}{}{0pt}{}[\vspace{1pt}\titlerule]

\titlespacing{\section}{
    % left space:
    -1pt
}{
    % top space:
    0.3 cm
}{
    % bottom space:
    0.2 cm
} % section title spacing

\renewcommand\labelitemi{$\vcenter{\hbox{\small$\bullet$}}$} % custom bullet points
\newenvironment{highlights}{
    \begin{itemize}[
        topsep=0.10 cm,
        parsep=0.10 cm,
        partopsep=0pt,
        itemsep=0pt,
        leftmargin=0 cm + 10pt
    ]
}{
    \end{itemize}
} % new environment for highlights


\newenvironment{highlightsforbulletentries}{
    \begin{itemize}[
        topsep=0.10 cm,
        parsep=0.10 cm,
        partopsep=0pt,
        itemsep=0pt,
        leftmargin=10pt
    ]
}{
    \end{itemize}
} % new environment for highlights for bullet entries

\newenvironment{onecolentry}{
    \begin{adjustwidth}{
        0 cm + 0.00001 cm
    }{
        0 cm + 0.00001 cm
    }
}{
    \end{adjustwidth}
} % new environment for one column entries

\newenvironment{twocolentry}[2][]{
    \onecolentry
    \def\secondColumn{#2}
    \setcolumnwidth{\fill, 4.5 cm}
    \begin{paracol}{2}
}{
    \switchcolumn \raggedleft \secondColumn
    \end{paracol}
    \endonecolentry
} % new environment for two column entries

\newenvironment{threecolentry}[3][]{
    \onecolentry
    \def\thirdColumn{#3}
    \setcolumnwidth{, \fill, 4.5 cm}
    \begin{paracol}{3}
    {\raggedright #2} \switchcolumn
}{
    \switchcolumn \raggedleft \thirdColumn
    \end{paracol}
    \endonecolentry
} % new environment for three column entries

\newenvironment{header}{
    \setlength{\topsep}{0pt}\par\kern\topsep\centering\linespread{1.5}
}{
    \par\kern\topsep
} % new environment for the header

\newcommand{\placelastupdatedtext}{% \placetextbox{<horizontal pos>}{<vertical pos>}{<stuff>}
  \AddToShipoutPictureFG*{% Add <stuff> to current page foreground
    \put(
        \LenToUnit{\paperwidth-2 cm-0 cm+0.05cm},
        \LenToUnit{\paperheight-1.0 cm}
    ){\vtop{{\null}\makebox[0pt][c]{
        \small\color{gray}\textit{Last updated in September 2024}\hspace{\widthof{Last updated in September 2024}}
    }}}%
  }%
}%

% save the original href command in a new command:
\let\hrefWithoutArrow\href

% new command for external links:


\begin{document}
    \newcommand{\AND}{\unskip
        \cleaders\copy\ANDbox\hskip\wd\ANDbox
        \ignorespaces
    }
    \newsavebox\ANDbox
    \sbox\ANDbox{$|$}

    \begin{header}
        \fontsize{25 pt}{25 pt}\selectfont Ahmet Neçirvan Doğan

        \vspace{5 pt}

        \normalsize
        \mbox{Istanbul, Türkiye}%
        \kern 5.0 pt%
        \AND%
        \kern 5.0 pt%
        \mbox{\hrefWithoutArrow{mailto:youremail@yourdomain.com}{andogan22@ku.edu.tr}}%
        \kern 5.0 pt%
        \AND%
        \kern 5.0 pt%
        \mbox{\hrefWithoutArrow{tel:+90-541-999-99-99}{+90 531 508 66 18}}%
        \kern 5.0 pt%
        \AND%
        \kern 5.0 pt%
        \AND%
        \kern 5.0 pt%
        \mbox{\hrefWithoutArrow{https://linkedin.com/in/yourusername}{www.linkedin.com/in/ahmet-necirvan-dogan-77a6b321a}}%
        \kern 5.0 pt%
        \AND%
        \kern 5.0 pt%
        \mbox{\hrefWithoutArrow{https://github.com/yourusername}{github.com/ahmetnecirvandogan}}%
    \end{header}

    \vspace{5 pt - 0.3 cm}


    \section{Experience}
        \begin{twocolentry}{
            Jul 2025 – Aug 2025
        }
            \textbf{Software Engineer Intern}, SimTek - Simulation Technologies -- Ankara, Türkiye
        \end{twocolentry}
        \vspace{0.10 cm}
        \begin{onecolentry}
            \begin{highlights}
                \item  I collaborated with the engineering team on a graphics-intensive software project, developing features with C++ and OpenGL.
            \end{highlights}
        \end{onecolentry}
           \vspace{0.10 cm} 
        \begin{twocolentry}{
            Jun 2024 – Sep 2024
        }
            \textbf{The International College Program / Cast Member}, Walt Disney World -- USA\end{twocolentry}

        \vspace{0.10 cm}
        \begin{onecolentry}
            \begin{highlights}
                \item  Collaborating within a diverse team of international and domestic cast members to ensure the efficient and safe operation of a high-volume attraction, demonstrating strong teamwork and interpersonal skills in a multicultural setting.
            \end{highlights}
        \end{onecolentry}
                

    
    \section{Education}
        
        \begin{twocolentry}{
            Oct 2022 – Jun 2027
        }
            \textbf{Koç University} Bachelor of Science - Computer Engineering\end{twocolentry}
            \begin{onecolentry}
            \textbf Minor: Media And Visual Arts
            \end{onecolentry}
            \begin{onecolentry}
            \textbf Track: Artificial Intelligence
            \end{onecolentry}

    
    \section{Projects}



        
        \begin{twocolentry}{
            Apr 2025 - Jun 2025
        }
            \textbf{
Interactive World Building Game}\end{twocolentry}

        \vspace{0.10 cm}
        \begin{onecolentry}
            \begin{highlights}
                \item An interactive 3D simulation that allows users to construct and manipulate a dynamic virtual world. 
                \item Built with OpenGL and C++, the simulation features a modular grid system, realistic lighting (using the Phong model), dynamic shadows, and a day-night mode/cycle. 
                \item The system emphasizes user-driven creativity, offering an immersive experience through real-time rendering techniques.
            \end{highlights}
        \end{onecolentry}

        \vspace{0.2 cm}

                \begin{twocolentry}{
            May 2025 - Jun 2025
        }
            \textbf{Sphere Renderer}\end{twocolentry}

        \vspace{0.10 cm}
        \begin{onecolentry}
            \begin{highlights}
                \item Developed a C++/OpenGL application featuring a physically simulated sphere with dynamic lighting and texturing. 

                \item Implemented core graphics techniques including Phong/Gouraud shading, custom PPM texture loading, multiple display modes (wireframe, shaded, textured), and planar shadow projection. 
                \item User interaction includes control over lighting, materials, camera zoom, and display toggles.
            \end{highlights}
        \end{onecolentry}

        \vspace{0.2 cm}

        \begin{twocolentry}{
            Apr 2025 - May 2025
        }
            \textbf{Rubik's Cube}\end{twocolentry}

        \vspace{0.10 cm}
        \begin{onecolentry}
            \begin{highlights}
                \item Built an interactive 3D Rubik’s Cube using modern OpenGL and GLSL. 
                \item Phong Lighting Technique for realistic rendering. 
                \item Modeled 27 sub-cubes with accurate face coloring and implemented smooth face rotations through transformation matrices.
                \item Developed a color-based picking system for selecting faces and included a help interface for user guidance. 
                \item The project emphasized real-time rendering, matrix manipulation, and shader-based graphics.
            \end{highlights}
        \end{onecolentry}

        
        \vspace{0.2 cm}

        \begin{twocolentry}{
            Mar 2025 - Apr 2025
        }
            \textbf{Bouncing Object Simulation}
            \end{twocolentry}

        \vspace{0.10 cm}
        \begin{onecolentry}
            \begin{highlights}
                \item Developed an interactive 3D application in shader-based OpenGL using GLFW, GLEW, GLM libraries with C/C++, simulating a horizontally bouncing object (cube, sphere, or bunny) with basic physics. 
                \item Enabled real-time object switching, color toggling, wireframe mode, and background rendering. Implemented event-driven interaction and transformations, emphasizing OpenGL pipeline understanding, shader programming, and fundamental physics simulation.
            \end{highlights}
        \end{onecolentry}

        \vspace{0.2 cm}

        \begin{twocolentry}{
            Dec 2023 - Jan 2024
        }
            \textbf{NBA Simulation Game}
            \end{twocolentry}

        \vspace{0.10 cm}
        \begin{onecolentry}
            \begin{highlights}
                \item Developed a comprehensive NBA simulation application in Java, employing OOP principles (Inheritance, Polymorphism), Swing for GUI development, and the Collections Framework for data management; implemented core functionalities including account creation, player drafting, season/playoff simulation, and file-based logging/stats processing.     
            \end{highlights}
        \end{onecolentry}


    
    \section{Technologies}



        
        \begin{onecolentry}
            \textbf{
            Languages:} C++, C, Java, Python, Scheme, MATLAB  \end{onecolentry}

        \vspace{0.2 cm}

        \begin{onecolentry}
            \textbf{
            Technologies:} OpenGL, JavaFX, XCode, Blender, Microsoft Office Programs
        \end{onecolentry}


    
    \section{Certifications}



        
        \begin{onecolentry}
            \textbf{
            Graph Search, Shortest Paths, and Data Structures} Stanford University  \end{onecolentry}

        \vspace{0.2 cm}

        \begin{onecolentry}
            \textbf{
            Computer Graphics with Modern OpenGL and C++} Udemy
        \end{onecolentry}

        \vspace{0.2 cm}

        \begin{onecolentry}
            \textbf{
            Disney Internships and Programs} The Walt Disney Company
        \end{onecolentry}

                \vspace{0.2 cm}

        \begin{onecolentry}
            \textbf{
            MATLAB Certified} Koç University
        \end{onecolentry}


    

\end{document}
